% Paso 1 (Alan). Sección 0 del documento.
\section{Sistema del Proyecto Fin de Curso}

\subsection{Roles del equipo}
% Paso 1 a) y b): rol, nombre completo, cédula y correo institucional de cada integrante.
% La guía tiene 5 roles para 3 integrantes: el grupo decide cómo se reparten.

\begin{center}
\small
\begin{tabularx}{\textwidth}{|>{\raggedright\arraybackslash}p{3.5cm}|L|>{\raggedright\arraybackslash}p{2.3cm}|>{\raggedright\arraybackslash}p{4.2cm}|}
\hline
\rowcolor{verdeuteq}\textcolor{white}{\textbf{Rol}} & \textcolor{white}{\textbf{Nombre completo}} & \textcolor{white}{\textbf{Cédula}} & \textcolor{white}{\textbf{Correo institucional}} \\ \hline
 & & & \\ \hline
 & & & \\ \hline
 & & & \\ \hline
\end{tabularx}
\end{center}

\subsection{Descripción del sistema}
% Paso 1 c): máximo 150 palabras. Nombre, problema que resuelve, usuarios principales y contexto organizacional.

UteqFit-AI es un sistema para el gimnasio de la UTEQ, que es gratuito. Se supone que actualmente el registro de los usuarios se hace de forma manual, lo que dificultaría controlar quién ocupa los cupos y quién dejó de asistir. Lo usarían los estudiantes, los encargados del gimnasio y Bienestar Universitario. Además, el sistema usaría inteligencia artificial para sugerir la rutina de cada usuario, y esa rutina la aprobaría un entrenador antes de asignarla, si el gimnasio cuenta con uno.

\subsection{Stakeholders y matriz poder/interés}
% Paso 1 d): los 5 stakeholders principales clasificados con la matriz poder/interés del PMBOK 7.a ed.

Los stakeholders se clasificaron con la matriz poder/interés del PMBOK~\cite{pmbok7}.

\begin{center}
\small
\begin{tabularx}{\textwidth}{|>{\raggedright\arraybackslash}p{3.5cm}|>{\centering\arraybackslash}p{1.6cm}|>{\centering\arraybackslash}p{1.6cm}|L|}
\hline
\rowcolor{verdeuteq}\textcolor{white}{\textbf{Stakeholder}} & \textcolor{white}{\textbf{Poder}} & \textcolor{white}{\textbf{Interés}} & \textcolor{white}{\textbf{Estrategia}} \\ \hline
Estudiantes & Bajo & Alto & Mantener informados \\ \hline
Responsable del gimnasio & Alto & Alto & Gestionar de cerca \\ \hline
Encargado de turno & Bajo & Alto & Mantener informado \\ \hline
Bienestar Universitario & Alto & Bajo & Mantener satisfecho \\ \hline
Entrenador (si existe) & Bajo & Alto & Mantener informado \\ \hline
\end{tabularx}
\end{center}

\subsection{Tipo de sistema}
% Paso 1 e): a qué tipo de sistema pertenece (según la secc. 2.5) y por qué.
